% =====================================================================
%  When Do China's Critical-Mineral Export Controls Bind?
%  Evidence from Graphite, Gallium-Germanium, and Antimony
%  Course: Introduction to China's Foreign Trade and Investment (ZJU)
%  All numerical results are reproduced by 04_scripts/analysis_thesis.py.
%  Upload flat to Overleaf: main.tex, t1-t8.tex, references.bib
% =====================================================================
\documentclass[12pt,a4paper]{article}

\usepackage[utf8]{inputenc}
\usepackage[T1]{fontenc}
\usepackage{lmodern}
\usepackage[margin=2.5cm]{geometry}
\usepackage{setspace}
\usepackage{amsmath,amssymb}
\usepackage{booktabs}
\usepackage{tabularx}
\usepackage{array}
\usepackage{float}
\usepackage{caption}
\usepackage[round,authoryear,sort]{natbib}
\usepackage[hidelinks]{hyperref}

\onehalfspacing
\setlength{\parindent}{1.5em}
\captionsetup{font=small,labelfont=bf,justification=justified}

\title{\bfseries When Do China's Critical-Mineral Export Controls Bind?\\[4pt]
\large\normalfont Evidence from Graphite, Gallium--Germanium, and Antimony}
\author{Mario D. Martinez\\[2pt]
\small Student ID: 22501254\\
\small Master in Chinese Economy --- Zhejiang University\\
\small Course: Introduction to China's Foreign Trade and Investment}
\date{\today}

\begin{document}
\maketitle

\begin{abstract}
\noindent
Between 2023 and 2024 China placed gallium, germanium, graphite, and antimony
under export-licensing rules. Because a license is not an export ban, whether
these controls actually reduced imports from China is an empirical question. I
build a monthly importer--product--source panel from UN Comtrade for
2019--2026 (53 reporting countries, 10 products) and estimate a Poisson
pseudo-maximum-likelihood source-share triple difference: a binding control
should lower China-source imports of a treated product relative to imports from
the rest of the world and relative to untreated products. The effects are
selective. Antimony oxides show a clean, persistent relative decline of about
81\%, concentrated among importers that were exposed to China and already had
other suppliers. Antimony metal and graphite point the same way but are less
clean; antimony ores are too thinly traded to read; and the aggregate
gallium--germanium trade codes show no decline. China's controls bind---but
only for some products and some importers, so diversification policy should be
chosen product by product rather than read off the legal announcement.
\end{abstract}

\medskip
\noindent\textbf{Keywords:} China; export controls; critical minerals; foreign
trade policy; antimony; graphite; gallium; germanium; PPML; supply chains.

\newpage
\tableofcontents
\newpage

% ============================ 1. INTRODUCTION ========================
\section{Introduction}\label{sec:intro}

\subsection{Motivation}

China mines and, more importantly, processes most of the minerals that sit at
the base of chips, batteries, optics, and defense goods
\citep{iea_2025_global_critical_minerals_outlook}. From July 2023 to September
2024 it placed four of them---gallium, germanium, graphite, and antimony---under
export-licensing requirements, citing the 2020 Export Control Law and national
security \citep{prc_2020_export_control_law}. The Ministry of Commerce (MOFCOM)
and the General Administration of Customs (GAC) issued the rules; the measures
arrived inside the same technology rivalry in which the United States had
already restricted semiconductor exports to China
\citep{bown-2020-howunitedstates, bown2024semiconductors}.

The key fact for an importer is that these are \emph{licensing} regimes, not
automatic export bans. MOFCOM publicly stated that compliant applications for
gallium and germanium, for graphite, and for antimony were being approved or
reviewed \citep{mofcom_2023_gallium_germanium_approvals,
mofcom_2024_graphite_approvals, mofcom_2024_antimony_license_review}. A license
can bind tightly, bind loosely, or work mainly as a signal. Which of these
happened is a question about trade flows, not something that follows from the
announcement.

This is also a new object for the study of economic statecraft. Most empirical
work on weaponized interdependence examines the \emph{demand} side, where the
United States restricts access to chips, software, or dollar clearing
\citep{bown-2020-howunitedstates, bown2024semiconductors}. Mineral export
controls are the \emph{supply} side: they restrict physical materials whose
production and processing are geographically concentrated and take years to
replicate elsewhere. The same legal instrument---an export license issued on
national-security grounds---should therefore bind very differently across
products, depending on China's share, on how cleanly the controlled item maps to
a trade code, and on whether importers can find other suppliers. That
heterogeneity is the object of this paper.

\subsection{Research question}

This paper asks: \emph{when do China's licensing-based critical-mineral export
controls reduce importers' reliance on China, and how do product scope and
importer adjustment capacity shape the answer?} The question is narrower than
``did the controls work?''. I look for trade-flow patterns consistent with a
binding, China-specific licensing wedge. I do not estimate enforcement
intensity, license-denial rates, inventories, coercive success, welfare, or
strategic intent---none of which the data can see.

\subsection{Contribution and main findings}

The contribution is an empirical, product-level test of measures that are
usually discussed only qualitatively. Using a source-share triple-difference
design on importer--product--source--month data, I find selective effects. The
central case is antimony oxides (Harmonized System code HS~282580): China-source
imports fall about 81\% relative to rest-of-world (ROW) imports and to untreated
products, the pre-trend is clean, and the decline deepens over the following
year. The effect is concentrated among high switching-potential importers---those
exposed to China who also had other suppliers---where the relative fall is about
87\%, while low-potential importers show no decline. Antimony metal (HS~8110) is
directionally consistent but has poor pre-trends; antimony ores (HS~261710) are
too thinly traded to interpret; natural graphite (HS~2504) is suggestive only;
and the aggregate gallium--germanium codes (HS~811292, HS~282560) show no
detectable decline. These nulls are informative: they reject the idea that every
controlled product mechanically collapses.

I am deliberately cautious about statistics. Treatment is assigned at the
product--month level, so with one treated and three untreated products in each
comparison, exact product-level inference cannot reach a $p$-value below 0.25.
Reporter-clustered $p$-values are reported as conditional diagnostics, not as
definitive assignment tests \citep{roth2022pretest}.

\subsection{Paper structure}

Section~2 reviews the literature and states the gap. Section~3 describes the
data. Section~4 sets out a simple licensing-wedge framework. Section~5 gives the
empirical strategy. Section~6 reports the results. Section~7 draws narrow policy
implications, and Section~8 concludes.

% ======================= 2. LITERATURE REVIEW ========================
\section{Literature Review}\label{sec:lit}

\subsection{Economic statecraft and chokepoints}

The conceptual anchor is the weaponized-interdependence research program.
\citet{farrellnewman-2019-weaponizedinterdependencehow} argue that global
economic networks are not neutral: states that sit at central nodes can turn
that position into coercive capacity. Their later book extends the argument to
the financial and digital infrastructure of the modern world economy
\citep{farrell2023underground_empire}. The mechanism is general---network
concentration plus jurisdictional control---but the empirical work it inspired
has focused almost entirely on the \emph{demand} side, where the United States
and its partners restrict access to technology and finance. The U.S.
semiconductor controls are the canonical case: \citet{bown-2020-howunitedstates}
shows how tariffs, entity-list actions, and export restrictions pulled the chip
industry into the trade war, and \citet{bown2024semiconductors} frame the same
conflict as a clash of industrial policies in which export controls are a core
tool. China's mineral measures are the supply-side mirror: they restrict
physical materials rather than designs or software, and they rest on geological
and processing concentration that is far slower to reverse. Recent work on
quantifying economic coercion stresses that the size of such effects must be
measured, not assumed \citep{eggeretal-2024-analyzingeffectseconomic}.

\subsection{Critical-mineral supply chains}

A separate literature explains why the same legal instrument should bind
unevenly across minerals. Supply is concentrated. The International Energy
Agency places antimony, gallium, germanium, and natural graphite among the most
concentrated critical minerals \citep{iea_2025_global_critical_minerals_outlook},
and the U.S.\ Geological Survey records China's dominant share of each
\citep{usgs_2026_antimony, usgs_2026_gallium, usgs_2026_germanium,
usgs_2026_graphite_natural}. Supply is also rigid in a specific way: gallium and
germanium are byproduct metals whose output tracks alumina and zinc, not their
own prices, so they respond sluggishly to a licensing shock
\citep{frenzel2017byproduct_metals}. Antimony adds a different rigidity, because
its processing, recycling, and substitution options are narrow
\citep{vandenbrink2022antimony_resilience, xu2025antimony_material_flow}, while
graphite's risk runs through the battery-anode processing chain rather than
through mining alone \citep{park2025graphite_supply_chain}. Two implications
follow for this paper. First, mineral export controls are structurally unlike
tariffs, financial sanctions, or technology bans, because adjustment requires
new mines and refineries on multi-year horizons. Second, heterogeneous effects
across products are theoretically expected, not surprising: the same license
should bind most where China's share is high, where the controlled item maps
cleanly onto a trade code, and where importers lack alternative sources.

One precedent sits in the background. The 2010 China--Japan rare-earth episode is
the clearest earlier case of a mineral being used as a foreign-policy instrument,
and it is often cited as evidence that mineral concentration translates into
coercive capacity. But it was a single event aimed at one country and lasting
weeks, whereas the 2023--2024 measures are open-ended licensing regimes that
apply to all importers and span several minerals. The precedent motivates the
question; it cannot answer it. Contemporary commentary on the 2023--2024
controls has been largely qualitative and forward-looking, and the supply-chain
literature describes vulnerability rather than realized trade responses.

\subsection{Gap}

What is missing is systematic, product-level evidence on whether and where these
specific Chinese controls actually shifted sourcing. Existing commentary on the
2023--2024 measures is qualitative or forward-looking; the concentration and
byproduct literature describes vulnerability but not realized trade responses;
and the one clear precedent, the 2010 rare-earth episode, is a single event
rather than a comparison across products. This paper fills that gap with a
quasi-experimental trade-flow test that treats the same licensing instrument
across several minerals and asks where it binds.

% ============================ 3. DATA ================================
\section{Data}\label{sec:data}

\subsection{UN Comtrade source-split panel}\label{sec:comtrade}

The main panel is built from monthly bilateral imports reported to United
Nations (UN) Comtrade, from January 2019 through March 2026 (87 months) for 53
reporting countries and ten products. For each reporter--product--month I keep
only the canonical aggregate record (partner2 code 0, customs code C00,
mode-of-transport and mode-of-shipment codes 0). This matters because the raw
files also contain world totals, bilateral partners, customs breakdowns, and
transport-mode breakdowns, and summing these double-counts trade. The balanced
country--product panel has $53\times10\times87=46{,}110$ rows; the
China/ROW source panel has 92,220 rows.

For each reporter--product--month I split imports into a China source (partner
156) and a ROW source (all other partners). The outcome is monthly trade value
in U.S.\ dollars; net weight is incomplete for several products and is used only
as a descriptive check. The three untreated control products---alumina
(HS~2818), manganese articles (HS~8111), and other inorganic acids (HS~2811)---are
large, regularly traded, and unaffected by any 2023--2024 licensing measure,
which is what a donor pool needs.

\subsection{Treatment chronology and product scope}\label{sec:chronology}

Table~\ref{tab:chronology} gives the three country-neutral licensing episodes
that define treatment. For the empirical work I date the post-period at the
\emph{expected import-effect month}, roughly one month after exporter
implementation, to allow a shipping lag. The three antimony headings are kept
separate because they sit at different value-chain stages (ores, metal, oxides)
and differ sharply in how regularly they are traded; pooling them would mix very
different responses.

A recurring issue is that the legal product scope is narrower than the
six-digit HS codes available in trade data, so the mapping from law to data is
imperfect. The gallium and germanium controls illustrate the extreme case: in
Comtrade, unwrought gallium and germanium are bundled with indium, vanadium, and
niobium under HS~811292, only two of which are controlled, and germanium dioxide
is bundled with zirconium dioxide under HS~282560. The controlled items are
therefore a fraction of trade in their headings, which mechanically pushes any
negative effect toward zero. The antimony oxide heading (HS~282580) is better,
because the controlled high-purity oxide dominates the heading even though the
legal scope is purity-thresholded. Graphite splits cleanly into a natural
heading (HS~2504, a good match for the controlled flake items) and an
artificial-graphite heading (HS~3801) that is only partially treated; I use
HS~2504 as the treated graphite product and treat HS~3801 as descriptive, never
as a placebo. Table~\ref{tab:products} summarizes the mapping and each
product's role.

\input{t1_chronology}
\begin{table}[H]
\centering
\caption{Product mapping and role in the analysis. ``Control'' products were
not subject to a Chinese export-licensing measure during the estimation window.}
\label{tab:products}
\small
\begin{tabular}{lll}
\toprule
HS-6 & Product & Role in the paper \\
\midrule
2504    & Natural graphite                              & Treated; suggestive \\
811292  & Ga/Ge/In/V/Nb unwrought bundle               & Treated; null / measurement-sensitive \\
282560  & Germanium dioxide / zirconium dioxide         & Treated; null / measurement-sensitive \\
8110    & Antimony metal and articles                   & Treated; consistent (pretrends weak) \\
261710  & Antimony ores and concentrates                & Treated; descriptive only \\
282580  & Antimony oxides                               & Treated; main result \\
\addlinespace
2818    & Alumina                                       & Control \\
8111    & Manganese articles                            & Control \\
2811    & Other inorganic acids                         & Control \\
\addlinespace
3801    & Artificial / colloidal graphite               & Descriptive; partially treated (not a placebo) \\
\bottomrule
\end{tabular}
\end{table}


\subsection{Coverage and regularity}\label{sec:coverage}

Table~\ref{tab:coverage} reports coverage for the six treated products. The
contrast between antimony oxides (HS~282580), with China value in 45\% of cells
and 43 active China importers before announcement, and antimony ores
(HS~261710), with China value in under 2\% of cells and only 12 active
importers, already signals that these two headings cannot play the same
evidentiary role. Sparse headings produce extreme coefficients that carry little
information.

\begin{table}[H]
\centering
\caption{Data coverage and trade regularity for the six treated products.
``China value share'' is the share of bilateral import value attributed to China
across the full panel (2019-01--2026-03). ``Positive China cells'' is the share
of importer-product-source-month cells with strictly positive China-source
imports in the pre-announcement window. ``Active China importers'' is the number
of reporters with at least one positive China-source flow before announcement.
Source: \texttt{comtrade\_thesis\_data\_audit.csv} and
\texttt{licensing\_wedge\_regularity\_diagnostics.csv}.}
\label{tab:coverage}
\footnotesize
\begin{tabular}{llccc}
\toprule
HS-6 & Product & China value share & Positive China cells & Active China importers \\
\midrule
2504   & Natural graphite                & 0.559 & 0.589 & 51 \\
811292 & Ga/Ge/In/V/Nb bundle           & 0.305 & 0.310 & 44 \\
282560 & GeO$_2$/ZrO$_2$                 & 0.457 & 0.477 & 48 \\
8110   & Antimony metal                  & 0.267 & 0.292 & 44 \\
261710 & Antimony ores/concentrates      & 0.016 & 0.019 & 12 \\
282580 & Antimony oxides                 & 0.419 & 0.449 & 43 \\
\bottomrule
\end{tabular}

\smallskip
\footnotesize Control products 2818, 8111, and 2811 have China value shares of
0.778, 0.577, and 0.862 respectively, confirming regular, well-populated
baseline trade required for a donor pool.
\end{table}


\subsection{Switching-potential measures}\label{sec:switchingdata}

For each importer--product pair I measure pre-announcement China exposure and
ROW access, and combine them into a switching-potential index
$A_{ik}=E_{ik}\times O_{ik}$, where $E_{ik}$ is the pre-period China share of
the importer's imports of product $k$ and $O_{ik}$ is the share of pre-period
months with any ROW import. $A_{ik}$ is high only when an importer was both
exposed to China and already buying from elsewhere. It is an observed proxy for
switching potential, not a true outside option. Reporters with no pre-trade are
excluded from the switching split but kept in the main analysis. The diagnostic
is revealing: for antimony oxides, 50 of 53 reporters have pre-trade and the
median $A_{ik}$ is 0.16; for antimony ores, 19 reporters have no pre-trade and
the median is zero.

\subsection{Eurostat CN8 case study}\label{sec:eurostatdata}

To probe the gallium--germanium aggregation problem I use Eurostat Comext at the
eight-digit Combined Nomenclature (CN8) level. Within the HS~811292 bundle I
separate unwrought gallium and germanium (controlled) from unwrought indium and
vanadium (not controlled), which serve as within-bundle comparators. The cleaned
panel is balanced for the current 27-member European Union (EU27): 27 countries,
four products, 72 months. It is sparse at the country--month level, so it is a
measurement case study, not a primary result.

% ====================== 4. THEORETICAL FRAMEWORK =====================
\section{Theoretical Framework}\label{sec:framework}

\subsection{A licensing-wedge model of sourcing}\label{sec:model}

The model is deliberately simple and reduced-form. Consider an importer $i$ that
buys product $k$ in month $t$ either from China ($M^{C}_{ikt}$) or from the rest
of the world ($M^{R}_{ikt}$). When a licensing control takes effect, buying from
China becomes more costly or less certain. I summarize all of that---delay, denial
risk, paperwork, uncertainty, reduced availability---in a single latent wedge
$W_{kt}\geq 0$. The wedge is the mechanism through which a legal rule becomes a
real restriction; it is not observed directly.

The wedge lowers the ratio of China-source to ROW-source imports. As a
conceptual expression, write the log of that ratio as
\begin{equation}
\ln\!\left(\frac{M^{C}_{ikt}}{M^{R}_{ikt}}\right)
= \alpha_{ik} + \lambda_{t} - W_{kt} + u_{ikt},
\label{eq:wedge}
\end{equation}
where $\alpha_{ik}$ is a fixed importer--product sourcing pattern, $\lambda_t$ is
a common time effect, and $u_{ikt}$ is noise. Equation~\eqref{eq:wedge} is
intuition only: bilateral trade has many zeros and is strongly heteroskedastic,
so estimation cannot take logs directly and uses PPML instead
\citep{santossilva2006gravity}. But it states the identifying idea cleanly---a
larger wedge means a smaller China-to-ROW ratio.

\subsection{Testable predictions}\label{sec:predictions}

The model gives three predictions that the data can test.

\paragraph{Sign.} If the control binds, then $W_{kt}>0$ after the policy, and the
China-source share of the treated product falls. In the PPML triple difference of
Section~5 this is a negative coefficient on the
treated$\times$China$\times$post interaction:
\begin{equation}
W_{kt}>0 \;\Longrightarrow\; \beta_{k}<0.
\label{eq:sign}
\end{equation}
The converse is not symmetric. A zero or positive $\beta_k$ does \emph{not} mean
the policy was irrelevant: routine license approvals, trade codes that bundle
controlled with uncontrolled items, rerouting, stockpiling, offsetting demand
shocks, or slow adjustment can all hide a wedge at the code level. So a clean
negative effect is a sufficient but not necessary sign of binding.

\paragraph{Magnitude.} Under PPML, $\beta_k$ has a direct proportional reading.
The implied change in the China-source share of product $k$ is
\begin{equation}
\text{percent change in China-source share} = 100\cdot\bigl(e^{\beta_k}-1\bigr).
\label{eq:magnitude}
\end{equation}
For the main case, antimony oxides, $\beta=-1.69$, so the model implies an
$81.5\%$ relative fall. The evaluation step is then to ask whether a fall of that
size is consistent with what actually happened to China, ROW, and total imports,
and whether it reflects source switching (China down, ROW and total up) or pure
compression (total down too).

\paragraph{Cross-section (who adjusts).} Not every importer should respond the
same way. Define an importer's pre-announcement China exposure and ROW access as
\begin{equation}
E_{ik}=\frac{\sum_{t<A_k} M^{C}_{ikt}}
              {\sum_{t<A_k} M^{C}_{ikt}+\sum_{t<A_k} M^{R}_{ikt}},
\qquad
O_{ik}=\frac{\#\{t<A_k: M^{R}_{ikt}>0\}}{\#\{t<A_k\}},
\qquad
A_{ik}=E_{ik}\times O_{ik},
\label{eq:switching}
\end{equation}
where $A_k$ is the announcement month. $A_{ik}$ is high only when an importer was
both exposed to China and already buying from elsewhere; it is an observed
switching-potential proxy, not a true outside option. The model predicts the
negative China-source effect should be concentrated among high-$A_{ik}$
importers---those who both needed to adjust and could.

A concrete example makes this tangible. Suppose an importer bought most of its
antimony oxide from China before the control and also bought small amounts from
European recyclers. When the license bites, the Chinese shipment is delayed or
denied, so the importer expands its European orders. Its total need for antimony
oxide may not change, but the \emph{share} bought from China falls---China down,
ROW up---which is exactly what the triple difference measures. An importer with
no European supplier simply loses supply and consumes less. The first importer
has high $A_{ik}$; the second has low $A_{ik}$. The cleanest relative decline
should sit in the first group.

\subsection{Scope: what the model is not}\label{sec:scope}

The model is reduced-form and the wedge is latent, so $\beta_k$ identifies only
the \emph{combined} effect of the license on the China-source share. It does not
identify the wedge level, the license-denial probability, enforcement intensity,
inventories, strategic intent, or welfare. Those would require a structural model
of licensing and supplier behavior that the data and the small number of treated
products cannot support, and that the project rules set aside for a master's
thesis. The evaluation below is therefore at the level of sign, magnitude, and
cross-section---a disciplined test of whether the reduced-form model organizes
the evidence, not a structural calibration.

% ======================= 5. EMPIRICAL STRATEGY =======================
\section{Empirical Strategy}\label{sec:strategy}

\subsection{PPML source-share triple difference}\label{sec:ddd}

The main estimator is a Poisson pseudo-maximum-likelihood (PPML) triple
difference (DDD) on the importer--product--source--month panel
\citep{santossilva2006gravity}. PPML is used because bilateral trade has many
zeros and strong heteroskedasticity, which bias the log-linear alternative. For
importer $i$, product $p$, source $s$, and month $t$,
\begin{equation}
\mathbb{E}[M_{ipst}\mid X]=
\exp\!\left[\beta\,(\mathit{Treated}_{p}\times \mathit{China}_{s}\times \mathit{Post}_{kt})
+\Gamma_{ip}+\Gamma_{is}+\Gamma_{pt}+\Gamma_{st}+\Gamma_{ps}+\Gamma_{it}\right].
\label{eq:ppml}
\end{equation}
The six pairs of fixed effects absorb importer--product scale, importer--source
relationships, product--month global shocks, source--month common shocks, average
product--source differences, and importer--month demand. The coefficient of
interest is $\beta$ on the triple interaction. It asks: did the China-source
share of the treated product change after the control, relative to the same
source comparison for the three untreated products? For each treated product the
estimation stack contains that product and the three controls (HS~2818, HS~8111,
HS~2811), and $\mathit{Post}_{kt}$ turns on at the expected import-effect month
in Table~\ref{tab:chronology}. Because treatment effects are likely
heterogeneous, I use this high-dimensional fixed-effects triple difference
rather than a simple two-way-fixed-effects design
\citep{dechaisemartindhaultfoeuille-2020-twowayfixedeffects}.

\subsection{Event studies and the switching split}\label{sec:eventstudies}

I replace the static interaction with event-time interactions to check pre-trends
and dynamics. Good pre-period coefficients should be near zero, and post-period
effects should appear after the expected import-effect month and then persist or
deepen. To avoid a long, noisy string of monthly coefficients, adjacent months
are binned (for example, months 4--6 and 7--12), with a pre-period window as the
omitted baseline. Because pre-trend tests can mislead in both directions, I read
them as one diagnostic among several, not as a gate \citep{roth2022pretest}. For
the mechanism I split reporters with pre-trade into high- and low-$A_{ik}$
groups at the median and re-estimate Equation~\eqref{eq:ppml} in each group. If
source switching drives the result, the negative effect should sit in the
high-$A_{ik}$ group. I also report a raw before/after comparison of China, ROW,
and total imports to separate switching from compression.

\subsection{Inference}\label{sec:inference}

Inference is the binding constraint. Treatment is assigned at the
product--month level: every importer faces the same rule for a given product, so
the 53 reporters give outcome variation but not 53 independent policy shocks.
With one treated and three control products, the minimum exact
product-assignment $p$-value is $1/(1+3)=0.25$. I therefore report
reporter-clustered $p$-values only as conditional diagnostics and weight the
\emph{joint pattern}---sign, timing, persistence, and heterogeneity---more than
any single $p$-value. A descriptive exporter interrupted time series (ITS) with
heteroskedasticity- and autocorrelation-consistent (HAC) standard errors is used
only as a timing check, not as a causal design.

The practical consequence is important for reading the results. A large negative
estimate with a small reporter-clustered $p$-value is not, by itself, a
causally identified effect of the license; it is a conditional summary of one
product relative to three controls. The case for a binding wedge rests instead
on whether several independent signals line up---a clean pre-trend, a sharp
break at the expected import-effect month, persistence or deepening over the
post-period, and concentration of the effect in the importers the framework
predicts. Where those signals line up, as for antimony oxides, the reading is
strong; where they conflict, the reading is correspondingly weak.

% ============================ 6. RESULTS =============================
\section{Results}\label{sec:results}

\subsection{Product-level source-share effects}\label{sec:productlevel}

Table~\ref{tab:main_ddd} reports the static source-share triple difference for
each treated product. The pattern is heterogeneous, as the framework predicts.
The gallium--germanium bundle (HS~811292) and the germanium-dioxide heading
(HS~282560) show small positive estimates and no decline---consistent with
code-bundling and byproduct-supply noise. Natural graphite (HS~2504) is negative
but imprecise ($-30.4\%$, $p=0.107$) and its pre-trend fails. Antimony metal
(HS~8110) and antimony ores (HS~261710) show large negative point estimates
($-86.8\%$ and $-98.9\%$) but their pre-trends fail badly, so they cannot carry
a causal claim alone. The standout is antimony oxides (HS~282580): a relative
decline of $81.5\%$, a reporter-clustered $p$ of $0.002$, and a joint pre-trend
$p$ of $0.680$. The next subsections take these cases in order of evidentiary
strength.

The exporter side is consistent with this timing. A descriptive
interrupted-time-series check on China-reported exports shows an implementation
break for the products with a clean importer-side signal, which is reassuring
about timing but is not a causal design because it has no untreated
counterfactual. The main evidence therefore comes from the importer-side triple
difference, where the untreated products provide the comparison.

\begin{table}[H]
\centering
\caption{Main source-share PPML triple-difference results. The dependent
variable is monthly bilateral import value in USD. The reported coefficient is
the percent effect of the \textit{Treated}$\times$\textit{China}$\times$\textit{Post}
interaction, with standard errors clustered by reporter. ``Pretrend $p$'' is the
joint $p$-value of the pre-period event-time coefficients. Reporter-clustered
$p$-values are conditional diagnostics, not product-level assignment inference
(Section~\ref{sec:inference}). Source:
\texttt{thesis\_ppml\_results.csv} and \texttt{thesis\_event\_pretrend\_tests.csv}.}
\label{tab:main_ddd}
\footnotesize
\begin{tabular}{llrrrl}
\toprule
HS-6 & Product & Effect (\%) & Reporter $p$ & Pretrend $p$ & Assessment \\
\midrule
2504   & Natural graphite           & $-30.4$  & $0.107$  & $0.020$ & Suggestive \\
811292 & Ga/Ge/In/V/Nb bundle      & $+19.3$  & $0.365$  & $<0.001$ & No aggregate negative effect \\
282560 & GeO$_2$/ZrO$_2$           & $+17.5$  & $0.271$  & $0.260$ & No aggregate negative effect \\
8110   & Antimony metal             & $-86.8$  & $<0.001$ & $<0.001$ & Consistent; pretrends fail \\
261710 & Antimony ores/concentrates & $-98.9$  & $0.003$  & $<0.001$ & Descriptive only \\
282580 & Antimony oxides            & $-81.5$  & $0.002$  & $0.680$ & Central result \\
\bottomrule
\end{tabular}
\end{table}


\subsection{Main case: antimony oxides}\label{sec:antimonyoxide}

HS~282580 is the central case because it combines a large estimate with clean
pre-trends and a persistent path. The event-study coefficients in
Table~\ref{tab:eventstudy} tell the story simply. The pre-period coefficients are
small and insignificant---what a clean case should look like. After the expected
import-effect month the estimates turn sharply negative from $-55.4\%$ in month
one, then steepen to $-71.7\%$ over months 4--6, $-90.2\%$ over months 7--12, and
$-94.6\%$ from month 13 on. The wedge does not fade; it widens.

\begin{table}[H]
\centering
\caption{Event-study coefficients for the main case, HS 282580 antimony oxides,
from the source-share PPML triple difference with the pre-period window
m6-to-m4 as the omitted baseline. Months are relative to the expected
import-effect month (2024-10). Reporter-clustered $p$-values are conditional.
Source: \texttt{thesis\_event\_study\_results.csv}.}
\label{tab:eventstudy}
\footnotesize
\begin{tabular}{lrr}
\toprule
Event window & Effect (\%) & Reporter $p$ \\
\midrule
m13 or earlier   & $-4.1$  & $0.879$ \\
m12 to m7        & $+7.3$  & $0.745$ \\
m6 to m4 (base)  & \multicolumn{1}{c}{---} & --- \\
m2               & $+12.1$ & $0.581$ \\
m1               & $+7.6$  & $0.635$ \\
p0               & $-8.8$  & $0.722$ \\
p1               & $-55.4$ & $0.013$ \\
p2               & $-53.7$ & $0.107$ \\
p3               & $-58.2$ & $0.061$ \\
p4 to p6         & $-71.7$ & $0.029$ \\
p7 to p12        & $-90.2$ & $<0.001$ \\
p13 or later     & $-94.6$ & $<0.001$ \\
\bottomrule
\end{tabular}
\end{table}


What kind of adjustment is this? Table~\ref{tab:subcompression} compares mean
monthly imports before and after the expected effect month. For antimony oxides,
China-source imports fall $27.7\%$ while ROW imports rise $315.2\%$ and total
imports rise $135.7\%$. The relative China-source decline therefore comes with a
large rise in non-China sourcing, not with a collapse in overall availability.
This fits source switching or demand expansion better than pure compression; with
trade-value data alone I cannot separate the two, and I cannot tell genuine
non-China production from rerouting.

\begin{table}[H]
\centering
\caption{Substitution versus compression diagnostics: raw percent change in mean
monthly trade value from the pre-announcement window to the post
expected-import-effect window, by source. A negative China change combined with a
positive rest-of-world (ROW) and total change points to source switching or
demand expansion rather than pure import compression. Post-start months: 2024-01
for graphite, 2023-09 for Ga/Ge, 2024-10 for antimony. Source:
\texttt{licensing\_wedge\_substitution\_compression\_diagnostics.csv}.}
\label{tab:subcompression}
\footnotesize
\begin{tabular}{llrrr}
\toprule
HS-6 & Product & China change (\%) & ROW change (\%) & Total change (\%) \\
\midrule
2504   & Natural graphite           & $-21.1$ & $+13.4$  & $-6.4$ \\
811292 & Ga/Ge/In/V/Nb bundle      & $+65.5$ & $+30.8$  & $+40.9$ \\
282560 & GeO$_2$/ZrO$_2$           & $+19.6$ & $+1.5$   & $+8.1$ \\
8110   & Antimony metal             & $-49.1$ & $+314.9$ & $+219.7$ \\
261710 & Antimony ores/concentrates & $-99.1$ & $+125.2$ & $+118.2$ \\
282580 & Antimony oxides            & $-27.7$ & $+315.2$ & $+135.7$ \\
\bottomrule
\end{tabular}
\end{table}


\subsection{Importer switching-potential mechanism}\label{sec:mechanism}

Table~\ref{tab:switching} reports the switching-potential split. For antimony
oxides it is the strongest mechanism result in the paper. Among the 25
high-$A_{ik}$ importers---exposed to China and already buying elsewhere---the
relative decline is $87.4\%$ ($p<0.001$) with a clean split pre-trend ($p=0.354$).
Among the 25 low-$A_{ik}$ importers the estimate reverses to $+103.1\%$
($p=0.035$). This is exactly what source switching predicts: the effect sits
where importers both needed to adjust and could. Because the high-$A_{ik}$ group
also has the clean pre-trend while the low-$A_{ik}$ group shows no decline, the
result is a mechanism, not just an average negative coefficient. I describe this
case as supporting the source-switching mechanism.

\begin{table}[H]
\centering
\caption{Importer switching-potential split. Reporters with positive
pre-announcement trade are split at the median of $A_{ik}=E_{ik}\times O_{ik}$,
and the source-share PPML triple difference is re-estimated within each split.
Reporter-clustered $p$-values are conditional. Source:
\texttt{thesis\_switching\_potential\_split\_results.csv} and
\texttt{thesis\_switching\_potential\_split\_pretrends.csv}.}
\label{tab:switching}
\scriptsize
\begin{tabular}{llrrrrrrl}
\toprule
HS-6 & Product & \shortstack{High-$A$\\effect (\%)} & \shortstack{High-$A$\\$p$} & \shortstack{High-$A$\\pretrend $p$} & \shortstack{Low-$A$\\effect (\%)} & \shortstack{Low-$A$\\$p$} & \shortstack{Low-$A$\\pretrend $p$} & Interpretation \\
\midrule
2504   & Natural graphite           & $-34.1$ & $0.080$  & $0.089$ & $+34.5$ & $0.623$ & $0.006$ & Suggestive \\
811292 & Ga/Ge/In/V/Nb bundle      & $+9.2$  & $0.648$  & $<0.001$ & $+705.4$ & $0.002$ & $<0.001$ & No support \\
282560 & GeO$_2$/ZrO$_2$           & $+1.4$  & $0.905$  & $0.267$ & $+30.2$ & $0.739$ & $<0.001$ & No support \\
8110   & Antimony metal             & $-86.8$ & $<0.001$ & $<0.001$ & $-64.2$ & $0.243$ & $<0.001$ & Consistent, not clean \\
261710 & Antimony ores/conc.        & $-100.0$ & $<0.001$ & $<0.001$ & \multicolumn{1}{c}{large $+$} & $<0.001$ & $<0.001$ & Not interpretable \\
282580 & Antimony oxides            & $-87.4$ & $<0.001$ & $0.354$ & $+103.1$ & $0.035$ & $0.161$ & Strong support \\
\bottomrule
\end{tabular}
\end{table}


The split is also informative by failing to rescue the other products. Graphite
moves the right way in the high-$A_{ik}$ group ($-34.1\%$, $p=0.080$; pre-trend
$p=0.089$) but only mildly, so it stays suggestive. Antimony metal is
directionally consistent but both splits fail their pre-trends. Antimony ores are
not interpretable: the high-$A_{ik}$ estimate is mechanically near $-100\%$ and
the low-$A_{ik}$ estimate is explosively positive, with failed pre-trends
throughout---the right conclusion for a heading where 19 of 53 reporters have no
pre-trade. Aggregate gallium--germanium is not rescued either: HS~811292 and
HS~282560 show no negative high-$A_{ik}$ effect, and HS~282560 is a particularly
clean null because its high-$A_{ik}$ pre-trend is acceptable ($p=0.267$) yet the
estimate is essentially zero.
\subsection{Graphite, gallium--germanium, and antimony caveats}\label{sec:caveats}

Graphite (HS~2504) is consistent with a weaker or noisier wedge: the full-sample
estimate is negative but imprecise, the pre-trend fails, and the high-$A_{ik}$
split sharpens it only mildly. Several features plausibly weaken the signal: the
policy is framed as an adjustment of an older regime rather than a wholly new
one, the legal scope is narrower than the code, and battery-anode demand was
rising fast before December 2023 \citep{park2025graphite_supply_chain,
usgs_2026_graphite_natural}. Graphite is suggestive, not confirming.

The aggregate gallium--germanium headings are meaningful nulls rather than
failures. HS~811292 bundles the controlled metals with uncontrolled indium,
vanadium, and niobium; HS~282560 bundles germanium dioxide with zirconium
dioxide; and both metals are byproducts whose supply tracks the host metal
\citep{frenzel2017byproduct_metals, usgs_2026_gallium, usgs_2026_germanium}. A
negative licensing signal is therefore attenuated and confounded, and the
positive full-sample estimates are best read as ``not detected at this level of
aggregation,'' not as ``no effect on the controlled items.'' These nulls are
useful precisely because they rule out a uniform collapse story.

The antimony family shows why pooling is wrong. Antimony metal is consistent with
a broader disruption---the estimate is large and negative and the substitution
diagnostic shows China imports falling while ROW imports jump---but its
pre-trends fail in both splits, so it is consistent with the framework, not a
clean case.
Antimony ores are descriptive only: trade is so sparse that the $-98.9\%$
estimate is driven by a handful of cells. Keeping the three headings separate is
what makes it visible that the binding-wedge evidence is concentrated in oxides.

\subsection{Eurostat CN8 case study}\label{sec:eurostat}

Does finer disaggregation within HS~811292 recover a gallium or germanium effect?
Table~\ref{tab:eurostat} says: not clearly. The gallium-versus-indium/vanadium
estimate is negative ($-59.8\%$), the germanium estimate is positive ($+25.1\%$),
and the pooled estimate is negative ($-32.7\%$), but all are far from
conventional significance ($p=0.31$, $0.42$, $0.36$), and the minimum exact
product-assignment $p$ is $1/3$. CN8 disaggregation moves gallium the way the
framework predicts, but the panel is too sparse to carry the claim. I use it as a
measurement case study that argues for caution about code aggregation, not as a
gallium--germanium finding.

\begin{table}[H]
\centering
\caption{Eurostat CN8 case-study results for gallium and germanium within the
HS 811292 unwrought bundle, using indium and vanadium as within-bundle
comparators on the balanced current-EU27 panel. Dependent variable is monthly
import value in EUR. The minimum exact product-assignment $p$-value is $1/3$ in
each case-specific comparison. Source: \texttt{thesis\_ppml\_results.csv}.}
\label{tab:eurostat}
\footnotesize
\begin{tabular}{lrrr}
\toprule
Comparison & Effect (\%) & Reporter $p$ & Exact min.\ $p$ \\
\midrule
Ga vs In/V            & $-59.8$ & $0.307$ & $0.333$ \\
Ge vs In/V            & $+25.1$ & $0.419$ & $0.333$ \\
Pooled Ga/Ge vs In/V  & $-32.7$ & $0.355$ & $0.167$ \\
\bottomrule
\end{tabular}
\end{table}


\subsection{How well does the model fit the data?}\label{sec:reading}

The point of writing down the model in Section~4 is to be able to judge it
against the estimates. The reduced-form licensing-wedge model makes three
predictions---sign, magnitude, and cross-section---and the data speak to each.

\paragraph{Sign.} The model predicts $\beta_k<0$ where the wedge binds. This
holds clearly for antimony oxides ($-81.5\%$) and antimony metal ($-86.8\%$),
weakly for graphite ($-30.4\%$, insignificant), and fails for the two
gallium--germanium codes ($+19.3\%$, $+17.5\%$). Antimony ores show the right
sign but the estimate is too unstable to count. So the sign prediction is met
for the products the model expects to bind (clean code, regular trade) and fails
where trade-code bundling should attenuate it---which is itself consistent with
the model's caveat in Equation~\eqref{eq:sign}.

\paragraph{Magnitude.} For antimony oxides the model implies an $81.5\%$
relative fall in the China-source share (Equation~\ref{eq:magnitude}). The raw
before/after moves in Table~\ref{tab:subcompression} are consistent with a fall
of that order: China-source imports drop $27.7\%$ while ROW imports rise
$315.2\%$ and total imports rise $135.7\%$. A relative China-source decline of
roughly four-fifths is therefore not an artifact of the regression; it matches a
real rebalancing of the import mix. The same table also settles the
switching-versus-compression question the model leaves open: for antimony oxides
the pattern is unambiguously \emph{switching plus expansion}---China down, ROW
and total up---not compression of total availability.

\paragraph{Cross-section.} The model predicts the effect concentrates among
high-$A_{ik}$ importers. For antimony oxides this is the cleanest result in the
paper: high-$A_{ik}$ importers show an $87.4\%$ relative fall with a clean
pre-trend, while low-$A_{ik}$ importers show a positive estimate. The cross-section
prediction is also directionally right for graphite and antimony metal, and it
cleanly \emph{fails to rescue} the gallium--germanium codes, where even the
high-$A_{ik}$ group shows no decline. A model that could fit any pattern would be
uninformative; this one predicts the right cases and is silent or rejected on the
others.

\paragraph{Verdict.} The reduced-form model organizes the evidence well where it
is meant to---regularly traded, cleanly coded products exposed importers---and
its failures are themselves informative, pointing to trade-code bundling and
sparse trade rather than to a uniformly effective or uniformly toothless policy.
It is not a structural estimate of the wedge: $\beta_k$ identifies the combined
effect of the license on the China-source share, not the denial probability or
the welfare loss, which the data cannot separate.

% ====================== 7. POLICY IMPLICATIONS =======================
\section{Policy Implications}\label{sec:policy}

The evidence justifies only narrow claims. Four points are earned.

First, distinguish legal announcement from binding disruption. The
gallium--germanium codes show that a control can be legally in force yet produce
no detectable aggregate decline, while antimony oxides show a large and
persistent shift. Risk assessment should track trade flows, not only the policy
register, and should expect heterogeneous outcomes.

Second, make diversification product- and stage-specific. Antimony ores, metal,
and oxides behave very differently under the same announcement, so treating
``antimony'' as one commodity misallocates effort. The binding constraint and the
feasible substitute differ by value-chain stage.

Third, exposure matters only when alternative capacity exists. The switching
split shows the antimony-oxide response is concentrated among importers exposed
to China \emph{and} already buying elsewhere. Exposure alone does not measure
vulnerability, and access alone does not measure resilience; both must be
present for observed switching. Stockpiles cushion short friction, but persistent
restrictions call for processing capacity and supplier development, which take
longer to build \citep{vandenbrink2022antimony_resilience,
xu2025antimony_material_flow}.

Fourth, China's coercive capacity is real but not automatic, and a binding
wedge can also accelerate substitution away from China. The antimony-oxide
diagnostic shows ROW imports more than tripling after the control. I do not read
this as coercive success or failure---the data show no strategic intent or
welfare effects---but it means the same mechanism that reduces China-source
imports can, over time, erode the concentration that gives the control its force
\citep{iea_2025_global_critical_minerals_outlook}.

% ========================== 8. CONCLUSION ============================
\section{Conclusion}\label{sec:conclusion}

China's 2023--2024 controls on gallium, germanium, graphite, and antimony are
licensing instruments, not export bans, and whether a license binds is an
empirical question. Using a monthly importer--product--source panel for
2019--2026 and a PPML source-share triple difference, I find selective and
heterogeneous responses. Antimony oxides (HS~282580) are the clean case: a
roughly 81\% relative decline in China-source imports, a clean pre-trend, a path
that deepens for over a year, and a switching-potential split that isolates the
effect among exposed importers with other suppliers. Antimony metal is
consistent with the framework but pre-trend-contaminated; antimony ores are too
sparse to read; graphite is suggestive; and the aggregate gallium--germanium
headings show no detectable decline. Together these results are consistent with
a selective licensing-wedge reading and reject the idea that every controlled
product collapses.

Two limits should travel with these findings. The design identifies a relative
China-source effect, not an absolute shortage, and the small number of treated
products caps inferential precision, so reporter-clustered $p$-values are
conditional rather than definitive. The framework is also reduced-form: it reads
trade-flow patterns consistent with a binding wedge, not the structural
determinants of enforcement.

The lesson for China's foreign trade and investment is that export licensing is
now a routine trade-policy tool, that its effect must be measured rather than
inferred from the announcement, and that the relevant differences run along
product scope, value-chain stage, trade regularity, and importer adjustment
capacity. The findings also hint at an investment response: persistent controls
create incentives for alternative mining and refining abroad and for processing
investment inside China-governed chains, though tracing that is beyond this
paper. The natural next steps---exact-grade price comparisons and a larger
product universe once later 2025 measures take effect---are left for a master's
thesis.

% ========================== REFERENCES ===============================
\renewcommand{\refname}{References}
\bibliographystyle{plainnat}
\bibliography{references}

\end{document}
